% mathdown-default: the default look. Included by
% defaults/mathdown-default.yaml after core.tex. Copy and adapt it to make
% your own style.

% Pandoc marks Markdown lists without blank lines between items as "tight"
% and removes all space between their items. Normal LaTeX lists don't do
% that, so switch it off.
\renewcommand{\tightlist}{}

% With toc: true, start the text on a new page after the contents.
\let\mdtableofcontents\tableofcontents
\renewcommand{\tableofcontents}{\mdtableofcontents\clearpage}

% Theorem heads on their own line, like the bold line in the Markdown:
%   **Theorem (Name)**        →   Theorem 3 (Name)
%   For all a, b …                For all a, b …
\newtheoremstyle{plainbreak}{\topsep}{\topsep}{\itshape}{}{\bfseries}{}{\newline}{}
\newtheoremstyle{defnbreak}{\topsep}{\topsep}{\normalfont}{}{\bfseries}{}{\newline}{}
\newtheoremstyle{rembreak}{\topsep}{\topsep}{\normalfont}{}{\itshape}{}{\newline}{}

% Numbering: one counter shared by all theorem-like blocks. With
% numbersections on, it is prefixed by the section and restarts in each
% section (Theorem 2.3); without, it runs 1, 2, 3, ...
\theoremstyle{plainbreak}
\ifnum\value{secnumdepth}>0
  \newtheorem{mdtheorem}{Theorem}[section]
\else
  \newtheorem{mdtheorem}{Theorem}
\fi
\newtheorem{mdlemma}[mdtheorem]{Lemma}
\newtheorem{mdproposition}[mdtheorem]{Proposition}
\newtheorem{mdcorollary}[mdtheorem]{Corollary}
\newtheorem{mdconjecture}[mdtheorem]{Conjecture}
\newtheorem{mdclaim}[mdtheorem]{Claim}
\theoremstyle{defnbreak}
\newtheorem{mddefinition}[mdtheorem]{Definition}
\newtheorem{mdexample}[mdtheorem]{Example}
\newtheorem{mdaxiom}[mdtheorem]{Axiom}
\newtheorem{mdnotation}[mdtheorem]{Notation}
\theoremstyle{rembreak}
\newtheorem{mdremark}[mdtheorem]{Remark}
